\section{A supply agreement does not establish available output}
\label{sec:apple-component-availability}

Once a transition identifies the components a product needs, the next question
is whether they can arrive on the required date. An agreement to purchase them
does not answer that operating question. Apple's component-availability risk
says that supply agreements cover many components but gives no assurance that
they can be extended or renewed on similar terms, or at all. It also identifies
supplier financial difficulties, business failure and industry consolidation as
possible limits on supply \citep[Risk Factors, printed p.~8]{apple2024}.
The historical record must therefore distinguish the agreement, its effective
term and the supplier's ability to perform. Renewal and delivery need their
own evidence before a scenario can rely on them.

More suppliers do not eliminate every shortage. Apple describes significant
supply and pricing risks for single or limited sources, while warning that
even components available from several sources can face industry-wide shortages
and commodity-price fluctuations. Its example of past semiconductor demand
and supply shortages reports an adverse effect on obtaining components and
products on commercially reasonable terms. The warning that such disruptions
could recur preserves a future possibility; the passage supplies no duration,
quantity shortfall or loss with which to calibrate it. A scenario may specify a
common industry constraint, but it must state that assumption rather than
derive independence or a diversification benefit from the number of suppliers.

Custom components present a different constraint. A sole source may initially
have insufficient capacity while yields mature or manufacturing expands.
Alternatively, the supplier may give priority to common components. A yield
problem and a production-allocation decision can both delay the same product,
yet the possible responses differ. Apple separately reports past adverse
effects from delayed or constrained component supply and delayed shipments of
completed products by outsourcing partners. Component production, final
assembly and delivery remain separate activities; an improvement in one need
not establish completion of the others. The Note 12 supply discussion and the
outsourcing passage retain the same distinction between sources, capacity,
partner performance and control \citep[Note 12, printed p.~46; Risk Factors,
printed p.~8]{apple2024}.

For an explicitly assumed delay, compare restoring the current source with
switching to another. The disclosure says that consequences can depend on the
time needed to obtain sufficient quantities from either route. A named
alternative is therefore only a proposed response until component
compatibility, available quantities, lead times, price and contractual terms
are established. Existing purchase obligations do not identify those facts or
the release of the old commitment. The model should retain both responses
with their separate prerequisites and delivery dates, leaving an infeasible
or unresolved response visible. Only after the relevant events and accounting
facts are specified can the resulting purchases, settlements or inventory
movements generate journals. The filing does not select a route or assign
probabilities to the alternatives.
